% =====================================================================
% 04-reward-signals.tex
% =====================================================================
\section{The Reward Signal}\label{sec:reward}

The reward is the only place where human intent enters RL training, so
its design dominates both the quality and the failure modes of the
resulting model. We survey the four reward sources identified in our
taxonomy and the orthogonal outcome-versus-process distinction.

\subsection{Human Preference Rewards: RLHF}

RLHF learns a reward model $r_\phi$ from pairwise human comparisons,
typically with a Bradley--Terry likelihood, and then optimizes the policy
against $r_\phi$ under a KL penalty to the reference
model~\cite{christiano2017deeprl,stiennon2020summarize,ouyang2022instructgpt}.
This recipe aligned InstructGPT and was later scaled for Llama~2-Chat,
which added rejection sampling before PPO and iterative rounds of
feedback collection~\cite{touvron2023llama2}.  Human preference remains
the gold standard for subjective qualities such as helpfulness and tone,
but it is expensive, noisy across annotators, and fragile: safety
alignment can be cheaply undone by subsequent
fine-tuning~\cite{qi2024finetuning,lermen2024lora}.  Comprehensive
treatments of RLHF and its limitations are given
in~\cite{kaufmann2024survey,casper2023open,lambert2024rlhfbook}.

\subsection{AI Preference Rewards: RLAIF}

To escape the annotation bottleneck, Constitutional AI has an LLM
critique and revise its own outputs under a written set of principles,
then trains a preference model on these AI
judgments~\cite{bai2022constitutional}.  RLAIF generalizes this idea and
was shown to match RLHF on summarization and dialogue while being far
cheaper and more reproducible~\cite{lee2023rlaif}.  The trade-off is that
AI feedback inherits the judge model's biases, including verbosity,
sycophancy, and self-preference, so RLAIF and RLHF are often combined,
using humans to calibrate and AI to scale.

\subsection{Implicit / Data Preference: DPO and Variants}

DPO observes that the RLHF objective has a closed-form optimum, allowing
the reward model to be eliminated and the policy to be trained directly
on preference pairs via a classification loss~\cite{rafailov2023dpo}.
This makes alignment dramatically simpler and more stable, at the cost of
being \emph{offline}: it learns from a fixed dataset and cannot explore.
A family of refinements followed: IPO regularizes against DPO's tendency
to overfit deterministic preferences~\cite{azar2024ipo}; KTO replaces
paired data with unpaired, prospect-theoretic
signals~\cite{ethayarajh2024kto}; Token-level Direct Preference
Optimization (TDPO) moves supervision to the token level for finer
control~\cite{zeng2024tdpo}; and calibrated DPO restores
scale-comparability with explicit rewards~\cite{liu2024caldpo}.  A key
empirical finding is that preference fine-tuning benefits from
suboptimal, on-policy data~\cite{tajwar2024preference}, foreshadowing the
online methods of Section~\ref{sec:optimization}.

\subsection{Verifiable / Rule-Based Rewards: RLVR}

The most consequential recent shift is toward rewards computed by a
deterministic verifier rather than a learned model: a math answer is
either correct or not, a unit test passes or fails.  Such
\emph{verifiable} rewards are cheap, unhackable in the usual sense, and
perfectly reproducible.  DeepSeekMath introduced GRPO precisely to
exploit them at scale~\cite{shao2024deepseekmath}, and DeepSeek-R1 showed
that pure RL with verifiable rewards, even without an initial SFT stage,
induces emergent chain-of-thought, self-verification, and ``aha''
behaviors~\cite{guo2025deepseekr1}.  Verifiable rewards are now the
dominant paradigm for reasoning, math, and code, and are a major reason
RL has become a third scaling axis alongside parameters and
data~\cite{snell2024testtime,openai2024o1,muennighoff2025s1}.

\subsection{Granularity: Outcome vs.\ Process Rewards}

Cutting across all four sources is the question of \emph{when} reward is
delivered.  Outcome reward models (ORMs) score only the final answer,
providing a sparse signal that complicates credit assignment over long
reasoning chains.  Process reward models (PRMs) instead score each
intermediate step; Lightman \textit{et~al.}\ showed PRMs outperform ORMs
for verifying mathematical reasoning~\cite{lightman2023verify}.  The
obstacle is the cost of step-level labels, addressed by automated process
supervision via Monte-Carlo rollouts~\cite{wang2024omegaprm} and by
careful study of how to build PRMs robustly~\cite{luo2025prmlessons}.
Recent work questions whether the PRM advantage transfers beyond
structured, verifiable domains, leaving process supervision for
open-ended tasks an open problem.
